\documentclass[11pt,a4paper]{article}
\usepackage[T1]{fontenc}
\usepackage[utf8]{inputenc}
\usepackage{charter}
\usepackage[margin=2cm]{geometry}
\usepackage[hidelinks]{hyperref}
\setlength{\parindent}{0pt}
\setlength{\parskip}{0.6em}
\pagestyle{empty}

\begin{document}

\textbf{Osinachi Chukwujama}\\
Saint-Étienne, France\\
\href{mailto:osinachi.chukwujama@gmail.com}{osinachi.chukwujama@gmail.com}\\
+33 6 51 41 94 14

\begin{flushright}
Recruitment Team, Electro-Optical Tests (EOT)\\
HGH Systèmes Infrarouges\\
10 rue Maryse Bastié\\
91430 Igny

\vspace{0.8em}
Saint-Étienne, 21 October 2026
\end{flushright}

\vspace{0.5em}
\textbf{Subject: Application for the electro-optical engineering internship, blackbody setup methods}

Dear Hiring Team,

I am a first-year student in the Erasmus Mundus Joint Master iPSRS (Intelligent Photonics), currently at Université Jean Monnet in Saint-Étienne, and I am writing to apply for your electro-optical engineering internship on blackbody setup methods.

Blackbodies sit at the base of every reliable infrared measurement, and this internship appeals to me because it combines hands-on production work with the methodical improvement of procedures. That is exactly the kind of environment where I have done my best work.

During my graduate internship at Rite Foods, I carried out daily maintenance on industrial baking and packaging machines, keeping downtime at or below 20\%. I wrote daily maintenance reports to keep equipment and supplies traceable, and when a part had no immediate replacement, I designed and 3D printed one. This taught me to work carefully on equipment alongside technicians, and to value clear records as much as the repair itself. My First Class degree in Mechanical Engineering and my current optics coursework, covering physical and Fourier optics, give me the technical base to understand the instruments I would be preparing and calibrating. In a laboratory practical this semester, I generated diffraction gratings and Fresnel lenses in MATLAB, displayed them on a spatial light modulator, observed the resulting diffraction patterns, and computed the gratings' spatial frequencies from my measurements.

IT tools are second nature to me. Between 2023 and 2026, I worked on a series of enterprise and freelance software projects, most recently as technical lead on a large intranet platform. That steady exposure taught me to pick up new software quickly, follow and improve procedures, and document my work so others can rely on it. I built a pipeline that generates formal documents automatically from database records, I write a technical blog, and as a Microsoft Certified Trainer I have prepared training material for more than 300 students. Turning feedback from the bench into clearer setup instructions and tracking sheets is work I would genuinely enjoy.

In preparing this application, I studied how blackbodies serve as radiometric references for infrared systems: how Planck's law links source temperature to emitted radiance, why emissivity, temperature uniformity and stability matter for calibration, and why burn-in is used to catch early failures before delivery. I am eager to see these principles applied on your test benches and to learn the practical side from your technicians and engineers. I am rigorous, I learn quickly, and I work well in a team.

My M1 semester ends on 31 May 2027, so I am available for three months from early June 2027. Thank you for considering my application. I would welcome the opportunity to discuss it with you.

Yours sincerely,

Osinachi Chukwujama

\end{document}
